\documentclass{article}

\usepackage[utf8]{inputenc}     % pdfLaTeX
\usepackage[T1]{fontenc}
\usepackage[magyar]{babel}

\usepackage{amsthm}
\usepackage{amsmath}
\usepackage{amsfonts}

\newtheorem{definition}{Definíció}
\newtheorem{corollary}{Folyomány}
\newtheorem{lemma}{Lemma}
\newtheorem{theorem}{Tétel}

\title{Gráfelmélet}
\author{Rovenszky Robin}
\date{Legutóbb frissítve: 2026. Március 26.}

\begin{document}

\maketitle
\tableofcontents
\newpage

\section{Bevezető}

\section{Definíciók}
\begin{definition}[Irányítatlan gráf]
    Csúcsokból és élekből áll. Egy él két, nem feltétlenül különböző csúcsot köt össze.
\end{definition}

\begin{definition}[Hurokél]
    Olyan él, amely egy csúcsot önmagához köt.
\end{definition}

\begin{definition}
    Egy él akkor többszörös él, ha létezik legalább egy másik él amely ugyan azt a két csúcsot köti össze.
\end{definition}

\begin{definition}[Egyszerű gráf]
    Olyan gráf, melynek nincs hurokélje és többszörös élje.
\end{definition}

\begin{definition}[Összefüggő gráf]
    Olyan gráf, melynek bármely pontjából bármely másik pontjába vezet út.
\end{definition}

\begin{definition}[Komponens]
    Olyan részgráf, amely összefüggő és a csúcsok számára nézve maximális.
\end{definition}

\begin{definition}[Fokszám]
    Adott csúcs fokszáma megadja az adott csúcsot elhagyó élek számát. Adott gráf fokszáma megadja az gráfban lévő csúcsok fokszámainak összegét.
\end{definition}

\begin{definition}[Üres gráf]
    Olyan gráf, melynek fokszáma 0.
\end{definition}

\begin{definition}[Teljes gráf]
    Olyan gráf, mely a fokszámára való tekintettel maximális, tehát minden csúcs minden csúccsal össze van kötve.
\end{definition}

\begin{definition}[Út]
    Egymáshoz kapcsolódó élek egymásutánja. Se él, se csúcs nem ismétlődik benne.
\end{definition}

\begin{definition}[Vonal]
    Olyan út, amelyben a csúcsok ismétlődhetnek.
\end{definition}

\begin{definition}[Séta]
    Olyan vonal, amely megengedi az élek ismétlődését.
\end{definition}

\begin{definition}[Kör]
    Olyan út, melynek kezdő és végpontja megegyezik.
\end{definition}

\begin{definition}[Hamilton-út]
    A gráf összes csúcsát tartalmazó út.
\end{definition}

\begin{definition}[Hamilton-kör]
    A gráf összes csúcsát tartalmazó kör.
\end{definition}

\begin{definition}[Euler-vonal]
    A gráf összes élét tartalmazó vonal.
\end{definition}

\begin{definition}[Euler-kör]
    A gráf összes élét tartalmazó kör.
\end{definition}

\begin{lemma}
    Egy összefüggő gráf komponenseinek száma 1.
\end{lemma}
\begin{proof}
    
\end{proof}

\section{Fák}
\begin{definition}[Fa]
    Egyszerű, összefüggő, körmentes gráf.
\end{definition}

\end{document}